\documentclass{article}

% packages
\usepackage{amsmath, amsthm, thmtools, amsfonts, amssymb, luacode, catchfile, tikzducks, hyperref, ifthen}
\ifcsname c@kobocompile\endcsname
	\usepackage[a5paper, total={1072pt, 1448pt}, margin=10pt, includeheadfoot]{geometry} % set page margins
\else
	\usepackage[a4paper, margin=50pt, includeheadfoot]{geometry}
\fi
\usepackage[shortlabels]{enumitem}
\usepackage[skip=3pt, indent=0pt]{parskip}

% language
\usepackage[bidi=basic, layout=tabular, provide=*]{babel}
\ifcsname c@english\endcsname
	\babelprovide[main, import]{english}
\else
	\babelprovide[main, import]{hebrew}
	\babelprovide{rl}
\fi
%\babelfont{rm}{Libertinus Serif}
\babelfont{rm}[Renderer=Harfbuzz]{Libertinus Serif}
\babelfont{sf}{Libertinus Sans}
\babelfont{tt}{Libertinus Mono}

% style
\AddToHook{cmd/section/before}{\clearpage}	% Add line break before section
\linespread{1.3}
\setcounter{secnumdepth}{0}		% Remove default number tags from sections, this won't do well with theorems
\AtBeginDocument{\setlength{\belowdisplayskip}{3pt}}
\AtBeginDocument{\setlength{\abovedisplayskip}{3pt}}
\graphicspath{ {../images/} }

% operators
\DeclareMathOperator\cis{cis}
\DeclareMathOperator\Sp{Sp}
\DeclareMathOperator\tr{tr}
\DeclareMathOperator\im{Im}
\DeclareMathOperator\re{Re}
\DeclareMathOperator\diag{diag}
\DeclareMathOperator*\lowlim{\underline{lim}}
\DeclareMathOperator*\uplim{\overline{lim}}
\DeclareMathOperator\rng{rng}
\DeclareMathOperator\Sym{Sym}
\DeclareMathOperator\Arg{Arg}
\DeclareMathOperator\Log{Log}
\DeclareMathOperator\dom{dom}
\DeclareMathOperator\supp{Supp}
\DeclareMathOperator\var{Var}
\DeclareMathOperator\cov{Cov}

% commands
%\renewcommand\qedsymbol{\textbf{מש''ל}}
%\renewcommand\qedsymbol{\fbox{\emoji{lizard}}}
\newcommand{\Aa}[0]{\mathcal{A}}
\newcommand{\Bb}[0]{\mathcal{B}}
\newcommand{\CC}[0]{\mathbb{C}}
\newcommand{\Cc}[0]{\mathcal{C}}
\newcommand{\EE}[0]{\mathbb{E}}
\newcommand{\FF}[0]{\mathbb{F}}
\newcommand{\Ff}[0]{\mathcal{F}}
\newcommand{\Ii}[0]{\mathcal{I}}
\newcommand{\Gg}[0]{\mathcal{G}}
\newcommand{\Ll}[0]{\mathcal{L}}
\newcommand{\Mm}[0]{\mathcal{M}}
\newcommand{\NN}[0]{\mathbb{N}}
\newcommand{\Nn}[0]{\mathcal{N}}
\newcommand{\PP}[0]{\mathbb{P}}
\newcommand{\Pp}[0]{\mathcal{P}}
\newcommand{\QQ}[0]{\mathbb{Q}}
\newcommand{\RR}[0]{\mathbb{R}}
\newcommand{\Rr}[0]{\mathcal{R}}
\newcommand{\Ss}[0]{\mathcal{S}}
\newcommand{\TT}[0]{\mathbb{T}}
\newcommand{\Uu}[0]{\mathcal{U}}
\newcommand{\Vv}[0]{\mathcal{V}}
\newcommand{\Ww}[0]{\mathcal{W}}
\newcommand{\ZZ}[0]{\mathbb{Z}}
\newcommand{\acts}[0]{\circlearrowright}
\newcommand{\explain}[2] {
	\begin{flalign*}
		 && \text{#2} && \text{#1}
	\end{flalign*}
}
\newcommand{\maketitleprint}[0]{ \begin{center}
	%\begin{tikzpicture}[scale=3]
	%	\duck[graduate=gray!20!black, tassel=red!70!black]
	%\end{tikzpicture}	
	\includegraphics[width=6cm]{cover}
\end{center}
}

% theorem commands
\newtheoremstyle{c_remark}
	{}	% Space above
	{}	% Space below
	{}% Body font
	{}	% Indent amount
	{\bfseries}	% Theorem head font
	{}	% Punctuation after theorem head
	{.5em}	% Space after theorem head
	{\thmname{#1}\thmnumber{ #2}\thmnote{ \normalfont{\text{(#3)}}}}	% head content
\newtheoremstyle{c_definition}
	{3pt}	% Space above
	{3pt}	% Space below
	{}% Body font
	{}	% Indent amount
	{\bfseries}	% Theorem head font
	{}	% Punctuation after theorem head
	{.5em}	% Space after theorem head
	{\thmname{#1}\thmnumber{ #2}\thmnote{ \normalfont{\text{(#3)}}}}	% head content
\newtheoremstyle{c_plain}
	{3pt}	% Space above
	{3pt}	% Space below
	{\itshape}% Body font
	{}	% Indent amount
	{\bfseries}	% Theorem head font
	{}	% Punctuation after theorem head
	{.5em}	% Space after theorem head
	{\thmname{#1}\thmnumber{ #2}\thmnote{ \text{(#3)}}}	% head content

\ifcsname c@english\endcsname
	\theoremstyle{plain}
	\newtheorem{theorem}{Theorem}[section]
	\newtheorem{lemma}[theorem]{Lemma}
	\newtheorem{proposition}[theorem]{Proposition}
	\newtheorem*{proposition*}{Proposition}
	%\newtheorem{corollary}[theorem]{אין חלופה עברית}

	\theoremstyle{definition}
	\newtheorem{definition}[theorem]{Definition}
	\newtheorem*{definition*}{Definition}
	\newtheorem{example}{Example}[section]
	\newtheorem{exercise}{Exercise}[section]

	\theoremstyle{remark}
	\newtheorem*{remark}{Remark}
	\newtheorem*{solution}{Solution}
	\newtheorem{conclusion}[theorem]{Conclusion}
	\newtheorem{notation}[theorem]{Notation}
\else
	\theoremstyle{c_plain}
	\newtheorem{theorem}{משפט}[section]
	\newtheorem{lemma}[theorem]{למה}
	\newtheorem{proposition}[theorem]{טענה}
	\newtheorem*{proposition*}{טענה}
	%\newtheorem{corollary}[theorem]{אין חלופה עברית}

	\theoremstyle{c_definition}
	\newtheorem{definition}[theorem]{הגדרה}
	\newtheorem*{definition*}{הגדרה}
	\newtheorem{example}{דוגמה}[section]
	\newtheorem{exercise}{תרגיל}[section]

	\theoremstyle{c_remark}
	\newtheorem*{remark}{הערה}
	\newtheorem*{solution}{פתרון}
	\newtheorem{conclusion}[theorem]{מסקנה}
	\newtheorem{notation}[theorem]{סימון}
\fi

% Questions related commands
\newcounter{question}
\setcounter{question}{1}
\newcounter{sub_question}
\setcounter{sub_question}{1}

\ifcsname c@english\endcsname
	\newcommand{\question}[1][0]{
		\ifthenelse{#1 = 0}{}{\setcounter{question}{#1}}
		\section{Question \arabic{question}}
		\addtocounter{question}{1}
		\setcounter{sub_question}{1}
	}

	\newcommand{\subquestion}[1][0]{
		\ifthenelse{#1 = 0}{}{\setcounter{sub_question}{#1}}
		\subsection{Part \alph{sub_question}}
		\addtocounter{sub_question}{1}
	}
\else
	\newcommand{\question}[1][0]{
		\ifthenelse{#1 = 0}{}{\setcounter{question}{#1}}
		\section{שאלה \arabic{question}}
		\addtocounter{question}{1}
		\setcounter{sub_question}{1}
	}

	\newcommand{\subquestion}[1][0]{
		\ifthenelse{#1 = 0}{}{\setcounter{sub_question}{#1}}
		\subsection{סעיף \localecounter{letters.gershayim}{sub_question}}
		\addtocounter{sub_question}{1}
	}
\fi

% import lua and start of document
\directlua{common = require ('../common')}

\GetEnv{AUTHOR}

% headers
\author{\AUTHOR}
\date\today

\title{פתרון מטלה 08 --- תורת הקבוצות (80200)}

\DeclareMathOperator\dom{dom}
\DeclareMathOperator\add{Add}
\DeclareMathOperator\mult{Mult}
\begin{document}
\maketitle
\maketitleprint{}

\directlua{ Q_number = 3; }
\Question{}
תהי $r$ קבוצה. נוכיח כי $\{ x : \exists y (\langle x, y \rangle \in r )\}$ היא קבוצה וגם $\{ y : \exists x(\langle x, y \rangle \in r )\}$ היא קבוצה. \\*
נסיק שאם $f$ פונקציה אז $\dom(f), \rng(f)$ קבוצות.
\begin{proof}
	נגדיר תכונה $p(x) = \exists y (\langle x, y \rangle \in r)$, ונוכל לבחור $\cup \cup r$ ואז מאקסיומת ההפרדה נקבל כי קיימת הקבוצה הראשונה, על־ידי הגדרת תכונה דומה והפעלת איחוד שלוש פעמים נקבל כי גם השנייה קיימת. \\*
	כמובן כל פונקציה $f$ היא קבוצת זוגות סדורים, וההגדרה של $\dom(f), \rng(f)$ היא שקולה לקבוצה הראשונה והשנייה שמצאנו כי קיימות זה עתה, ולכן נסיק כי התחום והתמונה של פונקציה הן קבוצות תמיד.
\end{proof}

\Question{}
תהי $X$ קבוצה. נוכיח במדויק על־ידי ZF את הטענות הבאות.

\Subquestion{}
נוכיח כי קיימת קבוצת כל יחסי השקילות על $X$.
\begin{proof}
	אנו יודעים כי יחס שקילות הוא יחס עם קיום תכונות מוגדרות נוספות, ולכן נגדיר $p(x)$ תכונה של קיום תכונות יחס שקילות. \\*
	מצאנו בהרצאה כי קבוצת הזוגות הסדורים $X \times X$ קיימת, ולכן מאקסיומת ההפרדה נוכל לטעון כי גם $\{ E \in X \times X \mid p(E) \}$ קיימת, וזו למעשה קבוצת כל יחסי השקילות על $X$.
\end{proof}

\Subquestion{}
נוכיח כי אם $E$ יחס שקילות על $X$ ו־$D$ מחלקת שקילות ביחס $E$, אז $D$ קבוצה.
\begin{proof}
	מצאנו עכשיו כי $E$ קבוצה, נגדיר $e \in X$ איבר כלשהו, ונגדיר $p(x) = \exists y \in X (\langle e, x \rangle \in E)$. \\*
	נשתמש שוב באקסיומת ההפרדה ונקבל כמובן כי $\{ x \in E \mid p(x) \}$ --- מחלקת שקילות המושרית על־ידי $e$ --- היא קבוצה. \\*
	נשאר לנו לקחת איבר כלשהו $e \in D$ ונקבל כי היא אכן קבוצה.
\end{proof}

\Subquestion{}
נוכיח שאם $E$ יחס שקילות על $X$, אז $X / E$ קבוצה.
\begin{proof}
	בסעיף הקודם מצאנו כי אם $x \in X$ אז מחלקת השקילות המכילה את $x$ היא קבוצה. \\*
	אנו יודעים כי $\mathcal{P}(X)$ היא קבוצה, ונבחין כי $D$ מחלקת השקילות כך ש־$x \in D$ מקיימת $D \in \mathcal{P}(X)$. \\*
	לכן אם כן נוכל להסיק כי התכונה $p(x) = \forall y, y' \in x \implies \langle \langle y, y' \rangle \rangle \in E$ ניתנת להפרדה ונקבל כי קבוצת מחלקות השקילות היא אכן קבוצה.
\end{proof}

\Question{}
תהי $X$ קבוצה, נוכיח שקיימת קבוצה $Y$ כך ש־$X \subseteq Y$ ולכל $y \in Y$ מתקיים $y \subseteq Y$.
\begin{proof}
	נגדיר פונקציה $g$ על־ידי $g(0) = X$ ו־$g(n + 1) = \bigcup_{a \in g(n)} a$. \\*
	על־ידי אקסיומת החלפה ומשפט הרקורסיה נסיק כי הפונקציה $g$ אכן קיימת ויחידה, ומאקסיומת האיחוד נקבל כי $Y = \bigcup_n g(n)$ אף היא קבוצה. \\*
	נותר רק לבדוק שהתנאי אכן מתקיים, יהי $y \in Y$, לכן קיים $n$ עבורו $y \in g(n)$, ואנו יודעים כי $y \subseteq \{ y \} \subseteq g(n + 1)$, ולכן קיבלנו כי $y \subseteq Y$ במתבקש.
\end{proof}

\Question{}
\Subquestion{}
ניעזר בפונקציה $\add : \NN \times \NN \to \NN$ כדי לבנות פונקציה $\mult : \NN \times \NN \to \NN$ המתארת על פעולת הכפל על הטבעיים.

נבחין כי על הפונקציה לקיים את התכונות:
\[
	\mult(0, 0) = 0
	\qquad
	\mult(n + 1, m) = \mult(n, m) + m
	\qquad
	\mult(n, m + 1) = \mult(n, m) + n
\]
כאשר אנו רושמים $X + n$ כסימון לביטוי $\add(X, n)$ ו־$X + 1$ כסימון ל־$s(X)$.

הוכחת קיום ויחידות הפונקציה $\mult$ כפי שהגדרנו אותה זה עתה נוכל לעשות בתהליך זהה להוכחה שראינו בהרצאה, נקבע סדרת פונקציות $f_n(m + 1) = m_n(m) + n$ עם בסיס $f_n(0) = 0$ ונקבל כי הן קיימות ומגדירות את $\mult$.
בשלב השני נוכיח באינדוקציה כי לכל $n$, לכל בחירת $m$, מתקיים גם $f_{n + 1}(m) = f_n(m) + m$.

\end{document}
